Это дополнительная лабораторная работа по подготовке собственной
Linux-среды. \textbf{Выполните только один маршрут}, подходящий вашему
компьютеру, затем общее задание в терминале. При установке Debian в
виртуальной машине Windows или macOS остаётся запущенной; при установке
второй системой ОС выбирают во время загрузки компьютера.

Работа с Debian нужна для знакомства с файловой системой, компилятором и
запуском программ Linux. Это отдельная подготовительная работа, а не
третий обязательный вариант семинара №4. Если подходящая Linux-среда уже
есть, можете сразу перейти к общей терминальной части и описать в отчёте
существующую установку.

\hypertarget{ux43fux43eux440ux44fux434ux43eux43a-ux441ux434ux430ux447ux438-ux434ux43eux43fux43eux43bux43dux438ux442ux435ux43bux44cux43dux43eux439-ux43bux430ux431ux43eux440ux430ux442ux43eux440ux43dux43eux439-ux43fux43e-debian}{%
\section{Порядок сдачи дополнительной лабораторной по
Debian}\label{ux43fux43eux440ux44fux434ux43eux43a-ux441ux434ux430ux447ux438-ux434ux43eux43fux43eux43bux43dux438ux442ux435ux43bux44cux43dux43eux439-ux43bux430ux431ux43eux440ux430ux442ux43eux440ux43dux43eux439-ux43fux43e-debian}}

\begin{enumerate}
\def\labelenumi{\arabic{enumi}.}
\tightlist
\item
  \textbf{На занятии} выполните доступную часть работы, зафиксируйте
  команды и наблюдения и покажите их преподавателю. Если в процессе
  работы возникли вопросы, задайте их преподавателю.
\item
  \textbf{К следующей лабораторной} завершите выбранный маршрут и
  практику, подготовьте отчёт: выбранный способ установки и причина
  выбора, команды, ожидаемый и фактический результаты, трудности и их
  решение. Изучите использованные команды и понятия; подготовьтесь
  объяснить, как работают программа, применённые механизмы и устройства
  компьютера. \textbf{Отправьте преподавателю отчёт в формате PDF или
  HTML.} Сохраните исходник \texttt{hello.c} отдельно для демонстрации;
  не включайте в отчёт пароли и ключи восстановления.
\item
  \textbf{Во время следующей лабораторной} индивидуально защитите
  работу: покажите отправленный отчёт, рабочую систему и исходник,
  повторите сборку и запуск, ответьте на вопросы. Преподаватель сообщит,
  если нужно что-то исправить.
\item
  \textbf{На последнем занятии} защитите предыдущую лабораторную работу
  и сдайте оставшиеся задолженности. Преподаватель объявит текущую
  накопленную оценку за лабораторные работы.
\end{enumerate}

Образец оформления собственного отчёта ищите рядом, в дополнительном
материале «Пример отчёта: сумма двух чисел». Его наблюдения относятся к
другой программе --- замените их своими.

\hypertarget{ux432ux44bux431ux435ux440ux438ux442ux435-ux43cux430ux440ux448ux440ux443ux442-ux43fux43bux44eux441ux44b-ux438-ux43cux438ux43dux443ux441ux44b}{%
\section{Выберите маршрут: плюсы и
минусы}\label{ux432ux44bux431ux435ux440ux438ux442ux435-ux43cux430ux440ux448ux440ux443ux442-ux43fux43bux44eux441ux44b-ux438-ux43cux438ux43dux443ux441ux44b}}

\textbf{Хост} --- уже работающая Windows или macOS; \textbf{гость} ---
Debian внутри виртуальной машины. На Windows x86-64 рекомендую начать с
VirtualBox: его диск легко отличить от диска Windows. На Mac с Apple
Silicon проще начать с UTM. При установке второй системой Debian
работает непосредственно на оборудовании и выбор ОС происходит при
включении компьютера. Для установки рядом с macOS сначала проверьте
поддержку конкретной модели по актуальной документации команды Debian.

\begin{longtable}[]{@{}
  >{\raggedright\arraybackslash}p{(\columnwidth - 4\tabcolsep) * \real{0.1800}}
  >{\raggedright\arraybackslash}p{(\columnwidth - 4\tabcolsep) * \real{0.3900}}
  >{\raggedright\arraybackslash}p{(\columnwidth - 4\tabcolsep) * \real{0.3900}}@{}}
\toprule
\begin{minipage}[b]{\linewidth}\raggedright
Подход
\end{minipage} & \begin{minipage}[b]{\linewidth}\raggedright
Плюсы
\end{minipage} & \begin{minipage}[b]{\linewidth}\raggedright
Минусы и границы
\end{minipage} \\
\midrule
\endhead
Windows: ВМ VirtualBox & Windows и Debian доступны одновременно; диск ВМ
можно копировать и восстанавливать. & Требуются ОЗУ и место на диске
хоста; доступ к устройствам ограничен. \\
Windows: вторая ОС & Нативная скорость и прямой доступ к оборудованию. &
Переключение через перезагрузку; разметка и загрузка затрагивают
реальный диск. \\
Mac: ВМ UTM & Быстрая arm64-виртуализация; macOS остаётся доступной. &
Это не Linux x86-64: задания NASM/x86 требуют другой среды. \\
Mac: Debian рядом с macOS & Нативная работа ARM без накладных расходов
ВМ. & Нужна поддержка модели установщиком команды Debian; переключение
требует перезагрузки, нет снимков ВМ. \\
\bottomrule
\end{longtable}

При установке из \textbf{netinst ISO} нужна сеть: установщик скачивает
недостающие пакеты. Названия пунктов интерфейса могут немного отличаться
в разных выпусках. Скачивайте \emph{актуальный стабильный выпуск}, не
привязываясь к конкретному номеру ISO: на
\href{https://www.debian.org/distrib/netinst}{официальной странице
Debian netinst} выберите \textbf{amd64} для Windows-ПК или
\textbf{arm64} для ВМ на Mac Apple Silicon. Эти ISO не взаимозаменяемы.
Перед использованием сравните контрольную сумму SHA256 скачанного файла
с файлом \texttt{SHA256SUMS} в том же официальном каталоге (в Windows
можно применить \texttt{Get-FileHash}, в macOS ---
\texttt{shasum\ -a\ 256}). Для установки Debian второй системой на Mac
применяется отдельный установщик; его шаги приведены в маршруте Г.

\hypertarget{ux43cux430ux440ux448ux440ux443ux442-ux430-debian-ux432-virtualbox-ux43dux430-windows}{%
\section{Маршрут А. Debian в VirtualBox на
Windows}\label{ux43cux430ux440ux448ux440ux443ux442-ux430-debian-ux432-virtualbox-ux43dux430-windows}}

\begin{enumerate}
\def\labelenumi{\arabic{enumi}.}
\tightlist
\item
  Установите \href{https://www.virtualbox.org/wiki/Downloads}{VirtualBox
  с официального сайта} в Windows. Если VirtualBox сообщает о
  недоступности аппаратной виртуализации, проверьте, включены ли
  \textbf{Intel VT-x / AMD-V} в настройках компьютера. Это возможность
  \emph{хоста}, не «вложенная виртуализация» внутри ВМ. Не меняйте
  параметры системы наугад: при затруднении обратитесь к инструкции
  производителя.
\item
  Скачайте Debian \textbf{amd64 netinst}. В VirtualBox нажмите
  \textbf{New / Создать}, назовите машину \texttt{debian-course},
  выберите ISO и систему \textbf{Debian (64-bit)}. Если мастер
  предлагает автоматическую установку, включите \textbf{Skip Unattended
  Installation / Пропустить автоматическую установку}, чтобы пройти
  учебные шаги вручную.
\item
  Выделите примерно \textbf{2 процессорных ядра} и \textbf{4096 МБ ОЗУ}
  (при меньшем объёме памяти компьютера допустимы 2048 МБ и лёгкое
  графическое окружение). Не выделяйте гостю всю память хоста. Создайте
  \textbf{новый динамический виртуальный диск} размером не меньше
  \textbf{30 ГБ}; убедитесь, что выбран именно новый диск ВМ, а не
  существующий файл с вашими данными.
\item
  Запустите ВМ с ISO и выберите \textbf{Graphical install}. Задайте
  язык, раскладку, имя компьютера \texttt{debian-course} и обычного
  пользователя. Если установщик предлагает пароль отдельного
  пользователя \texttt{root}, оставьте его пустым: согласно документации
  Debian первый обычный пользователь тогда получает возможность
  применять \texttt{sudo}. Не используйте пароль из примеров курса.
\item
  На этапе разметки выбирайте «Использовать весь диск» \textbf{только
  после проверки}, что в списке указан \textbf{виртуальный диск ВМ} (его
  размер должен совпадать с созданным вами, около 30 ГБ). Подойдёт один
  основной раздел для начинающего пользователя. Перед окончательным
  подтверждением прочитайте перечень действий установщика.
\item
  При выборе программ оставьте \textbf{standard system utilities} и
  выберите рабочий стол Debian (например, GNOME; на слабой ВМ можно
  выбрать Xfce). Установите GRUB на \textbf{диск ВМ}. Дождитесь
  завершения, отключите ISO от виртуального привода, если установщик
  предлагает это сделать, и перезагрузите ВМ. Войдите под созданным
  пользователем. После первой успешной загрузки можно сделать снимок ВМ
  через меню VirtualBox \textbf{Snapshots}.
\end{enumerate}

Диск ВМ --- это файл на Windows-компьютере; выбор «весь диск» в этом
маршруте \textbf{не относится к диску Windows}. Если не уверены, какой
диск показывает установщик, вернитесь к настройкам ВМ до подтверждения
разметки.

\hypertarget{ux43cux430ux440ux448ux440ux443ux442-ux431-windows-x86-64-ux441-uefi-debian-ux432ux442ux43eux440ux43eux439-ux441ux438ux441ux442ux435ux43cux43eux439}{%
\section{Маршрут Б. Windows x86-64 с UEFI: Debian второй
системой}\label{ux43cux430ux440ux448ux440ux443ux442-ux431-windows-x86-64-ux441-uefi-debian-ux432ux442ux43eux440ux43eux439-ux441ux438ux441ux442ux435ux43cux43eux439}}

Это работа с \textbf{реальным диском компьютера}. Прежде чем запускать
установщик, сохраните резервную копию важных файлов Windows на другом
носителе. В Windows проверьте через \texttt{msinfo32}, что в поле «Режим
BIOS» указано \textbf{UEFI}, а через «Управление дисками» --- что
системный диск размечен как \textbf{GPT}. Если используется BitLocker /
«Шифрование устройства», сохраните ключ восстановления доступным вне
этого компьютера и при необходимости временно приостановите защиту по
инструкции Windows. Без ключа или уверенности в схеме дисков выбирайте
маршрут А.

\begin{enumerate}
\def\labelenumi{\arabic{enumi}.}
\tightlist
\item
  \textbf{Подготовьте место в Windows.} В «Управлении дисками» уменьшите
  основной том Windows через \textbf{Shrink Volume / Сжать том} и
  оставьте \textbf{не менее 40 ГБ нераспределённого пространства}. Не
  создавайте на этом месте новый том в Windows, не форматируйте раздел
  EFI и не удаляйте раздел восстановления. Возможности уменьшения тома
  зависят от состояния диска; если Windows не даёт достаточно места,
  остановитесь и выбирайте ВМ вместо рискованных действий.
\item
  \textbf{Создайте установочную флешку.} Скачайте Debian \textbf{amd64
  netinst ISO}. Запишите его на USB-накопитель, например через
  \href{https://rufus.ie/}{Rufus} в Windows (схема разделов
  \textbf{GPT}, назначение \textbf{UEFI}). Проверьте выбранный
  USB-накопитель: запись удалит его прежнее содержимое. Полностью
  выключите Windows перед загрузкой с флешки; не продолжайте установку,
  если не подготовлены копия данных и нераспределённое место.
\item
  \textbf{Загрузите установщик в UEFI.} Через меню загрузки компьютера
  выберите пункт с флешкой, помеченный \textbf{UEFI}; не выбирайте
  вариант Legacy/CSM. В установщике Debian выберите \textbf{Graphical
  install}, язык, сеть, учётную запись. Как и для ВМ, при пустом пароле
  \texttt{root} первый пользователь получает \texttt{sudo}.
\item
  \textbf{Разместите Debian только в свободном месте.} В разделе
  разметки выберите автоматическую разметку \textbf{свободного
  нераспределённого места} (\emph{guided --- use largest continuous free
  space}, если этот вариант показан). \textbf{Никогда не выбирайте
  «Использовать весь диск» в этом маршруте.} Если выбора свободного
  места нет, остановитесь: не подменяйте его выбором всего диска. До
  записи изменений убедитесь, что системные разделы Windows (NTFS),
  раздел восстановления и уже существующий \textbf{EFI System Partition}
  не помечены для удаления или форматирования. Если план разметки
  непонятен --- отмените установку и спросите преподавателя или
  администратора.
\item
  Выберите рабочий стол Debian и \textbf{standard system utilities}.
  Установите загрузчик GRUB в \textbf{режиме UEFI}; существующий
  EFI-раздел Windows можно использовать для загрузочных файлов
  \textbf{без форматирования}. После установки извлеките флешку и
  перезагрузите компьютер. Выберите Debian в меню загрузки. Убедитесь,
  что Windows по-прежнему загружается через меню GRUB или меню UEFI
  компьютера. Если после изменений Windows просит ключ BitLocker,
  используйте сохранённый ключ; после проверок возобновите
  приостановленную защиту.
\end{enumerate}

Загрузчик, названия дисков и состав разделов зависят от конкретного ПК.
Не продолжайте установку, если в итоговом списке есть форматирование
раздела Windows или существующего EFI-раздела. Подготовка свободного
места через «Управление дисками» описана в
\href{https://learn.microsoft.com/en-us/windows-server/storage/disk-management/shrink-a-basic-volume}{документации
Microsoft}; возможности разметки --- в
\href{https://www.debian.org/releases/stable/amd64/ch06s03.en.html}{руководстве
установщика Debian}.

\hypertarget{ux43cux430ux440ux448ux440ux443ux442-ux432-mac-ux441-apple-silicon-debian-arm64-ux432-utm}{%
\section{Маршрут В. Mac с Apple Silicon: Debian arm64 в
UTM}\label{ux43cux430ux440ux448ux440ux443ux442-ux432-mac-ux441-apple-silicon-debian-arm64-ux432-utm}}

\begin{enumerate}
\def\labelenumi{\arabic{enumi}.}
\tightlist
\item
  Установите \href{https://getutm.app/}{UTM} на Mac и скачайте Debian
  \textbf{arm64 netinst ISO} с официальной страницы Debian. Не выбирайте
  ISO \texttt{amd64}: UTM будет \emph{виртуализировать} ARM-процессор, а
  не эмулировать x86-64.
\item
  В UTM нажмите \textbf{+}, выберите \textbf{Virtualize → Linux} и
  подключите скачанный ISO. Выделите, если позволяет Mac, около
  \textbf{4096 МБ ОЗУ}, \textbf{2 ядра} и виртуальный диск не меньше
  \textbf{30 ГБ}. Проверьте, что выбран режим \textbf{Virtualize}, а не
  \textbf{Emulate}. Имена некоторых переключателей зависят от версии
  UTM.
\item
  Запустите ВМ. В установщике Debian выберите \textbf{Graphical install}
  (если графический режим недоступен, используйте \textbf{Install} с
  теми же шагами), задайте язык, сеть, обычного пользователя. Оставьте
  пароль отдельного \texttt{root} пустым, чтобы первому пользователю был
  доступен \texttt{sudo}.
\item
  Для разметки разрешается «Использовать весь диск» \textbf{только на
  виртуальном диске UTM}. Выберите рабочий стол Debian и
  \textbf{standard system utilities}, завершите установку загрузчика на
  виртуальный диск и перезагрузите ВМ. Извлеките ISO из виртуального
  привода, если снова появляется установочное меню, и войдите в Debian.
\end{enumerate}

Выбор архитектуры влияет на дальнейшие задания: \texttt{uname\ -m} ниже
покажет \textbf{\texttt{aarch64}} (название архитектуры Debian ---
\texttt{arm64}). Компилятор GCC соберёт обычную C-программу из этой
работы. Исходники NASM для x86-64 и проверки курса, зависящие от
\textbf{Linux x86-64}, нельзя запускать как обычные arm64-программы; для
них понадобится другая x86-64-среда. Это ограничение относится и к
Debian, установленному непосредственно на Mac. Интерфейс создания ВМ
описан в \href{https://docs.getutm.app/basics/basics/}{документации
UTM}.

\hypertarget{ux43cux430ux440ux448ux440ux443ux442-ux433-mac-apple-silicon-debian-ux432ux442ux43eux440ux43eux439-ux441ux438ux441ux442ux435ux43cux43eux439}{%
\section{Маршрут Г. Mac Apple Silicon: Debian второй
системой}\label{ux43cux430ux440ux448ux440ux443ux442-ux433-mac-apple-silicon-debian-ux432ux442ux43eux440ux43eux439-ux441ux438ux441ux442ux435ux43cux43eux439}}

Здесь Debian устанавливается \textbf{рядом с macOS на реальный диск
Mac}. Обычный образ \texttt{arm64\ netinst}, предназначенный выше для ВМ
UTM, \textbf{не является установщиком этого маршрута}: для загрузки на
Mac нужны специальные компоненты. Команда Debian Bananas документирует
отдельный установщик с необходимым ядром и загрузочными компонентами из
своего репозитория. Перед началом откройте
\href{https://wiki.debian.org/Teams/Bananas}{страницу команды Debian
Bananas}, перейдите к актуальной инструкции установки на ARM Mac и
проверьте в ней \textbf{точную модель своего компьютера}. Если она не
указана как поддерживаемая, выберите маршрут В с UTM; не переносите шаги
для другой модели на свой Mac.

\begin{enumerate}
\def\labelenumi{\arabic{enumi}.}
\item
  В macOS сделайте резервную копию важных файлов на отдельном носителе и
  проверьте, что у вас есть место для второй ОС (рекомендуется не менее
  \textbf{40 ГБ}) и доступ к интернету. Откройте «Об этом Mac» и
  запишите модель для сверки с документацией команды Debian. Зарядите
  ноутбук или подключите питание на время установки.
\item
  Загрузите macOS как обычно и откройте приложение \textbf{Терминал}
  macOS. Сверив адрес с актуальной инструкцией команды Debian, запустите
  её установщик \textbf{в macOS}, не внутри ВМ:

\begin{verbatim}
curl -L https://bananas.debian.net/install | sh
\end{verbatim}
\item
  Следуйте сообщениям установщика: выберите, сколько места освободить от
  раздела macOS, какое окружение Debian установить (например, GNOME), и
  имя для записи Debian в выборе загрузочного диска. Установщик сам
  подготавливает нужные разделы и загрузочные компоненты; \textbf{не
  удаляйте разделы macOS или восстановления через «Дисковую утилиту»},
  не применяйте здесь инструкции Windows по GPT/EFI и не записывайте
  netinst ISO на флешку. Прежде чем подтвердить изменение диска,
  проверьте предлагаемые размеры. Если установщик не поддерживает модель
  либо список изменений не совпадает с ожидаемым, остановитесь и
  выберите ВМ.
\item
  После завершения первого этапа установщик попросит \textbf{выключить}
  Mac. Когда он полностью выключится, нажмите и удерживайте кнопку
  питания до появления \textbf{Loading startup options / Параметры
  запуска}. Выберите новый загрузочный диск Debian, при необходимости
  подтвердите действие учётной записью владельца Mac и завершите
  инструкции установщика. Дождитесь входа в Debian.
\item
  Проверьте обе системы: повторно откройте параметры запуска удержанием
  кнопки питания при выключенном Mac и загрузите macOS; затем аналогично
  загрузите Debian. На дальнейших занятиях \texttt{uname\ -m} в Debian
  покажет \texttt{aarch64}, как и в ВМ UTM. Нативная установка не
  превращает Apple Silicon в x86-64 и не позволяет запускать
  x86-64-программы NASM без отдельной подходящей среды.
\end{enumerate}

Для этого маршрута руководствуйтесь
\href{https://wiki.debian.org/Teams/Bananas}{текущими шагами команды
Debian} и связанным с ними
\href{https://asahilinux.org/docs/sw/partitioning-cheatsheet/}{описанием
разделов Apple Silicon}. Не путайте экран выбора загрузочного диска Mac
с меню GRUB для Windows-ПК: путь переключения ОС здесь другой.

\hypertarget{ux43eux431ux449ux430ux44f-ux43fux440ux430ux43aux442ux438ux43aux430-ux43fux435ux440ux432ux44bux439-ux432ux445ux43eux434-ux438-ux442ux435ux440ux43cux438ux43dux430ux43b}{%
\section{Общая практика: первый вход и
терминал}\label{ux43eux431ux449ux430ux44f-ux43fux440ux430ux43aux442ux438ux43aux430-ux43fux435ux440ux432ux44bux439-ux432ux445ux43eux434-ux438-ux442ux435ux440ux43cux438ux43dux430ux43b}}

В Debian с графическим рабочим столом найдите приложение
\textbf{Терминал} или попробуйте \textbf{Ctrl+Alt+T} (сочетание зависит
от окружения). Если выбрали текстовую установку, войдите под именем
своего пользователя: приглашение оболочки уже находится перед вами. При
наборе пароля символы обычно не показываются. Не копируйте знак
\texttt{\$} из примеров приглашения, вводите только команды.

\begin{verbatim}
whoami
uname -m
pwd
ls -la
mkdir -p ~/aiya/debian-first-steps
cd ~/aiya/debian-first-steps
pwd
touch notes.txt
nano notes.txt
cat notes.txt
\end{verbatim}

\texttt{whoami} покажет вашу учётную запись, \texttt{uname\ -m} ---
архитектуру Debian (\texttt{x86\_64} на Windows-маршрутах,
\texttt{aarch64} на обоих Mac-маршрутах). \texttt{pwd} показывает
текущий каталог, \texttt{ls\ -la} --- имена файлов, включая скрытые;
\texttt{\textasciitilde{}} означает домашний каталог. \texttt{mkdir\ -p}
создаст каталоги, \texttt{cd} перейдёт в рабочий каталог, \texttt{touch}
создаст пустой файл. В nano запишите строку «Я работаю в Debian»,
сохраните \textbf{Ctrl+O}, \textbf{Enter}, выйдите \textbf{Ctrl+X}.
\texttt{cat\ notes.txt} должен вывести введённую строку. Если
\texttt{pwd} не заканчивается на \texttt{/aiya/debian-first-steps},
вернитесь командой
\texttt{cd\ \textasciitilde{}/aiya/debian-first-steps} перед
продолжением.

Теперь подготовьте программы для учебных заданий. Команды с
\texttt{sudo} могут попросить пароль вашего пользователя:

\begin{verbatim}
sudo apt update
sudo apt install build-essential gdb nasm binutils nano file
gcc --version
make --version
gdb --version
nasm -v
\end{verbatim}

\texttt{apt\ update} обновляет список доступных пакетов,
\texttt{apt\ install} устанавливает инструменты.
\texttt{build-essential} включает GCC, \texttt{make} и файлы,
необходимые для сборки C; \texttt{gdb} --- отладчик, \texttt{nasm} ---
ассемблер x86, \texttt{binutils} --- набор утилит (включая \texttt{ld}),
\texttt{file} определяет тип файла. На Debian arm64 наличие NASM
\textbf{не превращает} гостевую систему в x86-64. Сравните версии
программ с выводом --- точные числа зависят от даты установки. Если
\texttt{sudo} сообщает, что пользователь не имеет прав, проверьте, как
была создана учётная запись при установке; обратитесь за помощью вместо
изменения системных настроек наугад.

\hypertarget{ux441ux43eux431ux435ux440ux438ux442ux435-ux438-ux437ux430ux43fux443ux441ux442ux438ux442ux435-ux43fux435ux440ux432ux443ux44e-ux43fux440ux43eux433ux440ux430ux43cux43cux443-c}{%
\section{Соберите и запустите первую программу
C}\label{ux441ux43eux431ux435ux440ux438ux442ux435-ux438-ux437ux430ux43fux443ux441ux442ux438ux442ux435-ux43fux435ux440ux432ux443ux44e-ux43fux440ux43eux433ux440ux430ux43cux43cux443-c}}

Оставайтесь в рабочем каталоге. Введите \texttt{nano\ hello.c}, наберите
текст \textbf{вручную} и сохраните его. При просмотре в PDF обратите
внимание на точки с запятой и двойные кавычки:

\begin{verbatim}
#include <stdio.h>

int main(void)
{
    printf("Hello from C!\n");
    return 0;
}
\end{verbatim}

Выполните сборку и проверку:

\begin{verbatim}
gcc -std=c11 -Wall -Wextra hello.c -o hello_c
./hello_c
echo $?
file hello_c
\end{verbatim}

GCC получает текст \texttt{hello.c} и создаёт исполняемый файл
\texttt{hello\_c}; флаги \texttt{-Wall\ -Wextra} включают полезные
предупреждения, \texttt{-o} задаёт имя результата. \texttt{./hello\_c}
просит Linux запустить программу из текущего каталога: ожидается
\textbf{\texttt{Hello\ from\ C!}}. \texttt{echo\ \$?} сразу после её
запуска покажет \textbf{\texttt{0}} --- код успешного завершения из
\texttt{return\ 0;}. Важно выполнить \texttt{echo\ \$?} \textbf{до}
следующей команды: после \texttt{file\ hello\_c} будет показан уже код
завершения программы \texttt{file}. Утилита \texttt{file} распознает
исполняемый ELF-файл; детали вывода будут различаться между
архитектурами \texttt{x86-64} и \texttt{AArch64}.

Если вместо результата видите \texttt{command\ not\ found}, проверьте
имя команды и установку пакета; для локального \texttt{hello\_c} нужен
префикс \texttt{./}. Для \texttt{No\ such\ file\ or\ directory}
проверьте \texttt{pwd} и \texttt{ls\ -la}: возможно, вы не в том
каталоге или опечатались в имени. Если GCC сообщил о синтаксической
ошибке, прочитайте указанные файл и строку, проверьте кавычки и
\texttt{;}, исправьте исходник в nano и \textbf{пересоберите}, прежде
чем запускать программу снова. Результат \texttt{Permission\ denied}
проверяйте через \texttt{ls\ -l\ hello\_c}: файл после GCC обычно уже
исполняемый, поэтому сначала выясните причину, а не применяйте
\texttt{sudo} к своей программе.

\hypertarget{ux447ux442ux43e-ux43fux43eux43aux430ux437ux430ux442ux44c-ux432-ux43eux442ux447ux451ux442ux435}{%
\section{Что показать в
отчёте}\label{ux447ux442ux43e-ux43fux43eux43aux430ux437ux430ux442ux44c-ux432-ux43eux442ux447ux451ux442ux435}}

Укажите выбранный маршрут и обоснуйте его с помощью таблицы плюсов и
минусов. Зафиксируйте, как вы создали виртуальный диск \textbf{или}
подготовили место для второй ОС, что увидели при первой загрузке, и
результаты \texttt{whoami}, \texttt{uname\ -m}, \texttt{pwd},
\texttt{ls\ -la}, установки инструментов и сборки. Запишите ожидаемый и
наблюдаемый вывод \texttt{hello\_c}, код завершения и одну возникшую
трудность либо явно сообщите, что ошибок не было. Отправьте отчёт
преподавателю \textbf{в PDF или HTML}, а \texttt{hello.c} сохраните
отдельно, чтобы повторить сборку и запуск на защите.
